%! Diagonalizing a Matrix — Unit 6, Lesson 6.2 — Warm-Up (Answer Key)
\documentclass[10pt]{article}
\usepackage{linalg-article}
\usepackage{linalg-key}   % pulls in -boxes; do NOT also load -boxes
\newcommand{\TallMath}[1]{$\displaystyle #1\rule[-1.4em]{0pt}{3.2em}$}

\begin{document}

\pageheader{Unit 6, Lesson 6.2}{Warm-Up --- Answer Key}
\namedateperiod

\vspace{0.12in}

\begin{notesbox}{Getting Ready: The Three Moves We'll Use Today}
\small Yesterday you \emph{found} eigenvalues and eigenvectors. Today you line them up into matrices and
put them to work --- and these three moves are the whole setup. Answer quickly.

\vspace{4pt}
\textbf{1. Build a matrix from columns.} Put the vectors $\mathbf{u}=\begin{bsmallmatrix}1\\1\end{bsmallmatrix}$
and $\mathbf{w}=\begin{bsmallmatrix}1\\-1\end{bsmallmatrix}$ into the \emph{columns} of a $2\times2$ matrix $X$:
\[
  X=\begin{bmatrix}{\color{keyred}\mathbf{1}} & {\color{keyred}\mathbf{1}}\\[3pt]{\color{keyred}\mathbf{1}} & {\color{keyred}\mathbf{-1}}\end{bmatrix}.
\]

\vspace{4pt}
\textbf{2. The $2\times2$ inverse (Unit 2).} Recall
$\begin{bsmallmatrix} a & b \\ c & d \end{bsmallmatrix}^{-1}=\dfrac{1}{ad-bc}\begin{bsmallmatrix} d & -b \\ -c & a \end{bsmallmatrix}$.
Find the inverse of $M=\begin{bmatrix} 1 & 1 \\ 0 & 1 \end{bmatrix}$:
\[
  M^{-1}={\color{keyred}\mathbf{\begin{bmatrix} 1 & -1 \\ 0 & 1 \end{bmatrix}}}.
\]

\vspace{4pt}
\textbf{3. Scaling the columns.} Multiplying on the \emph{right} by a diagonal matrix scales each column.
Compute
\[
  \begin{bmatrix} 1 & 1 \\ 1 & -1 \end{bmatrix}\begin{bmatrix} 3 & 0 \\ 0 & 1 \end{bmatrix}
  ={\color{keyred}\mathbf{\begin{bmatrix} 3 & 1 \\ 3 & -1 \end{bmatrix}}}.
\]
\end{notesbox}

\begin{teachernote}
These three moves \emph{are} today's diagonalization setup, in order: \textbf{(1)} line the eigenvectors up
as the \emph{columns} of $X$; \textbf{(2)} invert $X$ with the Unit~2 formula (here $\det M=1$, so no
fractions); \textbf{(3)} see that right-multiplying by a diagonal matrix \emph{scales each column} by its
eigenvalue --- exactly what $X\Lambda$ does. Item~1 and item~3 use the very $X$ from today's worked example
$A=\begin{bsmallmatrix}2&1\\1&2\end{bsmallmatrix}$ (eigenvectors $\begin{bsmallmatrix}1\\1\end{bsmallmatrix},
\begin{bsmallmatrix}1\\-1\end{bsmallmatrix}$). Debrief quickly, then name the plan: ``today we combine these
into $A=X\Lambda X^{-1}$.''
\end{teachernote}

\end{document}
